\documentclass[oneside,12pt]{book}
\usepackage[utf8]{inputenc}
\usepackage{standalone}
\usepackage[margin=2.5cm]{geometry}
\usepackage{graphicx}
\usepackage{hyperref}
\usepackage{lmodern}
\usepackage{float}
\usepackage[spanish]{babel}
\usepackage{csquotes}
\usepackage[figurename=Ilustración]{caption}
\usepackage[numbers]{natbib}
\usepackage{ragged2e}
\usepackage{array}
\usepackage{multirow}
\usepackage{parskip}
\usepackage{listings}
\usepackage{epigraph}
\usepackage{lipsum} % Texto de relleno en plantilla

\usepackage{algorithm}
\usepackage{algpseudocode}

% Glosario de términos - BEGIN
\usepackage{glossaries}

\makeglossaries

\newglossaryentry{latex}
{
    name=latex,
    description={Is a markup language specially suited 
    for scientific documents}
}

\newglossaryentry{maths}
{
    name=mathematics,
    description={Mathematics is what mathematicians do}
}
% Glosario de términos - END

%\usepackage{fancyhdr}
\usepackage{geometry}
\usepackage{color}\usepackage{graphicx,xcolor}

\renewcommand{\lstlistingname}{Algoritmo}
%\renewcommand{\lstlistlistingnam%e}{Índice de~\lstlistingname s}

\floatname{algorithm}{Algoritmo}
\renewcommand{\listalgorithmname}{Índice de~\lstlistingname s}


\lstset{
        basicstyle=\fontsize{7}{11}\ttfamily,
        language=Java,
        captionpos=b,
}

\usepackage{titlesec}
\setcounter{secnumdepth}{3}
\titleformat{\chapter}[display]
{\normalfont\huge\bfseries}{\chaptertitlename\ \thechapter}{20pt}{\Huge}

\urlstyle{same}
\title{Diseño y desarrollo de un sistema empotrado basado en bus CAN con acceso inalámbrico para la monitorización de un vehículo automóvil}
\author{Gorka Eymard Santana Cabrera}
\date{Curso 2025 - 2026}
\renewcommand{\contentsname}{Contenido}
\renewcommand{\figurename}{Ilustración}
\usepackage{pifont}
\renewcommand\labelitemi{\ding{52}}

\begin{document}
%\maketittle
\input{Portada.tex}

\input{TFT04C.tex}

\newpage
\pagenumbering{roman}
{\Large{\textbf{Agradecimientos}}}
\input{Agradecimientos.tex}

\clearpage
{\setlength{\parskip}{8mm}
{\huge{\textbf{Resumen}}}

La estructura de este documento es orientativa. Las \textbf{partes obligatorias (motivación y objetivos, competencias, aportaciones, alineamiento ODS, desarrollo, conclusiones y bibliografía)} deben estar recogidas, aunque pueden reordenarse o redistribuirse en base a las características del trabajo.

    {EL presente Trabajo de Fin de Título(TFT) desarrolla un sistema empotrado basado en bus CAN que trabaja en la capa de aplicación de OBDII permitiendo la realización de diagnósis de vehículos. El sistema se ha implementado utilizando como vehículo de pruebas un SEAT Ibiza 6j con motor 1.4 BXW(\textit{sustituir codigo motor por ECU}).}

\newpage

{\huge{\textbf{Abstract}}}


    {\large{Ahora en inglés: 
    This Final Degree Project (TFT) develops an embedded system based on the CAN bus that operates at the OBD-II application layer, enabling vehicle diagnostics. The system has been implemented using a SEAT Ibiza 6J with a 1.4 BXW engine (\textit{replace engine code with ECU designation}) as the test vehicle.
 }

}}
\tableofcontents

%\addcontentsline{toc}{section}{Listings}
%\lstlistoflistings
\listofalgorithms 

\listoffigures

\listoftables


\clearpage
\pagenumbering{arabic}

\chapter{Introducción}
\input{Introduccion.tex}

Ejemplo para el glosario de términos (ver al final del documento):

This final report has been produced using \Gls{latex}, a tool specially suitable for technical documents that include \gls{maths}. 

\chapter{Estado actual y objetivos iniciales}
%\input{EstadoDelArte.tex}
Este es un CAPÍTULO \textbf{OBLIGATORIO}.

\section{Motivación y antecedentes}

Descripción de la motivación inicial y el estado actual en la temática concreta del trabajo.

\section{Objetivos}

Desglose de los \textbf{objetivos} inicialmente previstos en el TFT01. Opcionalmente, pueden adelantarse las desviaciones que hayan podido producirse con respecto a esos objetivos iniciales.

NOTA: deben evitarse expresiones con referencias directas a formularios del tipo ``en el TFT01 ...'', puesto que cuando la memoria sea consultada fuera del entorno de la EII no serán identificables. Deben sustituirse por ``en la propuesta inicial ...''.

\chapter{Aportaciones del trabajo}
\chaptermark{Competencias, aportaciones y ODS}
%\input{Competencias.tex}
Este es un CAPÍTULO \textbf{OBLIGATORIO}.

También pueden incluirse parte de estos contenidos en las conclusiones, en la descripción de los objetivos o en forma de anexos.

\section{Principales aportaciones}

Justificar qué es lo que este TFT \textbf{aporta} a nuestro entorno socio-económico, técnico o científico. Repercusión esperada.

\section{Competencias específicas}

Indicar, sólo para las \textbf{competencias específicas} relacionadas de forma más directa con el trabajo desarrollado, cómo se han cubierto con este TFT.

\section{Alineamiento con los objetivos de desarrollo sostenible}

\begin{table}[!htbp]
\caption{Grado de relación del TFT con los objetivos de desarrollo sostenible.}
\label{tab:ODS}
\centering
\begin{tabular}{|l|c|c|c|c|}
\cline{1-5}
  & \multicolumn{4}{c|}{Grado de relación} \\ \cline{2-5} 
 ODS & 0 & 1 & 2 & 3 \\
  & No procede & Bajo & Medio & Alto \\  
  \cline{1-5} 
1 Fin de la Pobreza & & & & \\ \cline{1-5}
2 Hambre cero & & & & \\
\cline{1-5}
3 Salud y Bienestar & & & & \\
\cline{1-5}
4 Educación de calidad & & & & \\
\cline{1-5}
5 Igualdad de género & & & & \\
\cline{1-5}
6 Agua limpia y saneamiento & & & & \\
\cline{1-5}
7 Energía Asequible y no contaminante & & & & \\
\cline{1-5}
8 Trabajo decente y crecimiento económico & & & & \\
\cline{1-5}
9 Industria, Innovación e Infraestructuras & & & & \\
\cline{1-5}
10 Reducción de las desigualdades & & & & \\
\cline{1-5}
11 Ciudades y comunidades sostenibles & & & & \\
\cline{1-5}
12 Producción y consumo sostenibles & & & & \\
\cline{1-5} 
13 Acción por el clima & & & & \\
\cline{1-5}
14 Vida submarina & & & & \\
\cline{1-5}
15 Vida de ecosistemas terrestres & & & & \\
\cline{1-5}
16 Paz, justicia e instituciones sólidas & & & & \\
\cline{1-5}
17 Alianzas para lograr objetivos & & & & \\
\cline{1-5}
\end{tabular}%
\end{table}

Justificar el alineamiento del TFT con los \textbf{ODS} con los que presente un mayor grado de relación, en función de lo indicado en la tabla~\ref{tab:ODS}.

\chapter{Desarrollo}
%\input{Desarrollo.tex}
Este es un CAPÍTULO \textbf{OBLIGATORIO}.

\section{Metodología}

Descripción de la \textbf{Metodología} aplicada a lo largo del trabajo y justificación de la elección.

\section{Fases de desarrollo}

Desarrollo del trabajo en sus distintas fases (por ejemplo, para una metodología en cascada: análisis, diseño, implementación, validación). Deberán comentarse especialmente las decisiones de diseño tomadas y la selección de las herramientas empleadas.

Si el trabajo incluye software desarrollado, deberán seleccionarse las secciones más relevantes del mismo y comentarlas en la memoria.

Ajuste a la planificación inicialmente prevista.

Modificación en los objetivos planteados.

\chapter{Resultados}
%\input{Resultados.tex}

Presentación de los resultados del trabajo. Repercusión.

\chapter{Conclusiones y trabajo futuro}
%\input{Conclusiones.tex}
Este es un CAPÍTULO \textbf{OBLIGATORIO}.

\section{Conclusiones}

Valoración de resultados, grado de consecución de los  objetivos.

\section{Trabajo futuro}

Líneas de trabajo futuro, aspectos pendientes, posibles extensiones.

\section{Uso de la IA}

Indicar el uso que se ha hecho de la IA tanto en la elaboración de este documento como en el desarrollo del TFG.

\vspace{1cm}

\textbf{EXTENSIÓN DE LA MEMORIA:}

\begin{itemize}
    \item GII, GIFM, DGII-ADE (parte de GII): entre 50 y 100 páginas
    \item GCID: entre 75 y 100 páginas
\end{itemize}

\chapter{Otros capítulos y anexos}

El documento debe terminar con las \textbf{referencias bibliográficas} (SECCIÓN \textbf{OBLIGATORIA}). Después pueden incluirse \textbf{anexos} de forma opcional.

\section{Ejemplos de capítulos opcionales}

Dependiendo del tipo de trabajo, se podrían incluir, en el orden que corresponda, capítulos adicionales como los siguientes:

\begin{description}
\item{REQUISITOS}
\item{NORMATIVA Y LEGISLACIÓN} incluir la legislación vigente que afecte al TFT (ley de protección de datos, leyes sobre seguridad, ...)
\item{ASPECTOS ECONÓMICOS Y TEMPORALES}
\item{DISEÑO}
\item{RESULTADOS EXPERIMENTALES}
\item{MANUAL  DE  USUARIO  Y  SOFTWARE}  deberán  incluirse  obligatoriamente  en  la  memoria  los extractos más relevantes   del código desarrollado. Siempre que sea posible, deberá   proporcionarse acceso a un repositorio software.
\end{description}


\bibliography{ref}
\bibliographystyle{apalike}
\clearpage

\printglossaries

\end{document}